\section{Introduction}
\subsection{Transverse-field Ising model}
The underlying Hamiltonian for our discussion is the one-dimensional transverse-field Ising model with an additional longitudinal component
\begin{equation}
    \label{eq:TFI}
    \hat{H} = -J \sum_{i} \hat{\sigma}^x_i \hat{\sigma}^x_{i+1} - g \sum_i \hat{\sigma}^z_i - h \sum_i \sigma_i^x
\end{equation}
with Pauli matrices $\hat{\sigma}_i^j$, transverse-field strength $g$, longitudinal-field strength $h$, and coupling strength $J$. ~\cite{labmanual}
For $h = 0$, we recover the transverse-field Ising model, which we can solve analytically:
first, we set $J=1$ and perform a a Jordan-Wigner transformation, i.e.
mapping the spin operators to fermionic operators $\hat{c}_j$:
\begin{align*}
    \hat{\sigma}_j^z &= \mathbbm{1} - 2\hat{c}_j^\dagger \hat{c}_j \\
    \hat{\sigma}_j^- \coloneqq \frac{1}{2} \left( \hat{\sigma}_j^x - i \hat{\sigma}_j^y\right) 
    &= \left( - 1\right)^{\sum_{k < j} \hat{c}_k^\dagger \hat{c}_k} \cdot \hat{c}_j^\dagger \\ 
    \hat{\sigma}_j^+ \coloneqq \frac{1}{2} \left( \hat{\sigma}_j^x + i \hat{\sigma}_j^y\right) 
    &= \left( - 1\right)^{\sum_{k < j} \hat{c}_k^\dagger \hat{c}_k} \cdot \hat{c}_j.
\end{align*}
Following that, a discrete fourier transform of the form $\hat{c}_j = 1/\sqrt{N}\sum_k e^{-ikj} \hat{c}_k$ is performed
which yields the Hamiltonian in momentum representation 
\begin{equation*}
 \hat{H} = \sum_k [ 2 \cos k - 2 g ] \hat{c}_k^\dagger \hat{c_k} + \text{const}.
\end{equation*}
Ultimately, a Bogoliubov transformation diagonalizes the Hamiltonian $\hat{H} = 
\sum_k \varepsilon_k \hat{d}_k^\dagger \hat{d}_k$ and we obtain the quasiparticle
energy dispersion 
\begin{equation}
    \label{eq:dispersion}
    \varepsilon_k^\pm = \pm  2 \sqrt{(\cos k - g)^2 + \sin^2 k}
\end{equation}
This energy gap closes at $g = J = 1$ and $k = \pi$, indicating the presence
of a quantum phase transition at the critical point $g_\text{crit} = 1$. 
The system transitions from a ferromagnetic phase ($g < 1$) to a paramagnetic phase ($g > 1$). ~\cite{dutta2015,Coleman_2015}

When computing the spectral function later on, we will recover a shape similar of equation (\ref{eq:dispersion}), i.e., the $h=0$ case.

\subsection{Matrix product states (MPS)}
A generic many-body quantum state of a spin-$\frac{1}{2}$ system of size $N$ can be written as 
\begin{equation}
    \label{eq:psi}
    \ket{\psi} = \sum_{j_1, j_2, \dots, j_N} \psi_{j_1, j_2, \dots, j_N} \ket{j_1} \otimes \ket{j_2} \otimes \dots \otimes \ket{j_N},
\end{equation}
where $\psi_{j_1, j_2, \dots, j_N}$ is an order $N$ tensor. The idea of a MPS is to rewrite equation (\ref{eq:psi}) as a product of matrices

\begin{align}
    \ket{\psi} &= \sum_{j_1, \dots, j_N} \sum_{\alpha_2, \dots, \alpha_N} M^{[1]j_1}_{\alpha_1 \alpha_2} M^{[2]j_2}_{\alpha_2 \alpha_3} \dots M^{[N]j_N}_{\alpha_N \alpha_{N+1}} \ket{j_1, j_2, \dots, j_N} \notag \\
    &\equiv \sum_{j_1, \dots, j_N} M^{[1]j_1} M^{[2]j_2} \dots M^{[N]j_N} \ket{j_1, j_2, \dots, j_N}.
\end{align}
Each $M^{[n]j_n}$ is a $\chi_n \times \chi_n$ dimensional matrix, where $\chi_n$ is denoted as the bond dimension. 
The $\alpha_n$'s are referred to as "virtual" indices and the $j_n$'s are the so called "physical" indices.
The MPS representation can be obtain by successive reshaping and applying singular value decompositions (SVDs) to the initial coefficient tensor $\psi_{j_1, j_2, \dots, j_N}$.\cite{Hauschild2018}

This ansatz is very useful for numerical simulations. Unlike exact diagonalization methods whose computational cost usually grows exponentially with system size $N$, 
MPS algorithms typically scale polynomially in system size as $\mathcal{O}(\text{poly}(N))$.

\subsection{Time-Evolving Block Decimation (TEBD)}
When using the TEBD algorithm, we are interested in evaluating the time evolution of a quantum state
\begin{equation*}
    \ket{\psi(t)} = \hat{U}(t) \ket{\psi(0)}.
\end{equation*}
The time evolution can either be performed for real times $t$, or for imaginary times, i.e., 
performing a Wick rotation $\tau = i t$. ~\cite{Coleman_2015}
The latter yields the ground state of the system, given as 
\begin{equation*}
    \ket{\psi_0} = \lim_{\tau \to \infty} \frac{\hat{U}(\tau) \ket{\psi (0)}}{|| \hat{U}(\tau) \ket{\psi (0)}  ||}. \text{~\cite{labmanual}}
\end{equation*}
We consider a Hamiltonian with nearest-neighbor interactions and split it into even and odd indices 
\begin{equation*}
    \hat{H} = \sum_{n \text{ odd}} \hat{h}^{[n, n+1]} + \sum_{n \text{ even}} \hat{h}^{[n, n+1]},
\end{equation*}
divide the time evolution operator into small time slices $\delta t$, and  decompose the time-evolution operator by making use of the Suzuki-Trotter decomposition in linear order
\begin{equation}
    \hat{U}(\delta t) = \prod_{n \text{ odd}} \hat{U}^{[n, n+1]} (\delta t) \prod_{n \text{ even}} \hat{U}^{[n, n+1]} (\delta t) + \mathcal{O}(\delta t^2).
\end{equation}

This approximation makes it apparent that, due to the first-order nature, the Trotter error
accumulates over long simulation times.
Roughly speaking, the heart of the TEBD algorithm consists of the successive application of these two-site
unitary operators, which locally update our MPS, i.e. quantum state, with every time step. ~\cite{Hauschild2018}

\subsection{Dynamical Spin Structue factor}

The excitation spectrum of a quantum system is encoded in the spin structure factor $S(k, \omega)$, which is formally given by
the double Fourier transform of the real-time spin-spin correlation fucntion in space and time $C(t,j)$. For the numeric implementation,
the fourier transforms become discretized and we multiply by a Gaussian window function $G(t_n)$ to avoid Gibbs oscillations.
\begin{equation*}
    S(k, \omega) = \sum_{j=0}^{L-1} \sum_{t_n = 0}^N e^{i \omega \Delta t \cdot t_n - i k \cdot j} C (\Delta t \cdot t_n, j) \cdot G(t_n). \text{~\cite{labmanual}}
\end{equation*} 